\documentclass[12pt,a4paper]{article}

\usepackage[T2A]{fontenc}
\usepackage[utf8]{inputenc}
\usepackage[russian]{babel}
\usepackage{amsmath,amssymb,amsthm,mathtools}
\usepackage{geometry}
\usepackage{microtype}
\usepackage{enumitem}
\usepackage{xcolor}
\usepackage{tikz}
\usepackage{hyperref}

\geometry{left=25mm,right=20mm,top=22mm,bottom=24mm}
\hypersetup{
  colorlinks=true,
  linkcolor=blue!45!black,
  urlcolor=blue!55!black
}
\setlist{nosep}
\setlength{\parindent}{1.25em}
\setlength{\parskip}{0.15em}

\newtheorem{theorem}{Теорема}
\newtheorem{proposition}{Утверждение}
\newtheorem{lemma}{Лемма}
\theoremstyle{definition}
\newtheorem{definition}{Определение}
\theoremstyle{remark}
\newtheorem*{remark}{Замечание}

\DeclareMathOperator{\Sp}{Sp}
\DeclareMathOperator{\diag}{diag}
\DeclareMathOperator{\cond}{cond}
\newcommand{\CC}{\mathbb{C}}
\newcommand{\RR}{\mathbb{R}}
\newcommand{\ZZ}{\mathbb{Z}}
\newcommand{\eps}{\varepsilon}
\newcommand{\norm}[1]{\left\lVert #1\right\rVert}

\DeclareMathOperator{\ind}{ind}
\newcommand{\Conv}{\mathcal C}

\title{Операторы Винера--Хопфа и их факторизация\\Односторонняя обратимость}
\author{Конспект лекции}
\date{31 октября 2026 года}

\begin{document}
\maketitle

\section{Оператор Винера--Хопфа как сжатие свёртки}

Сохраним соглашения предыдущей лекции. На $X=\ell^1(\ZZ)$ задан
координатный проектор
\[
  (P_+x)_k=\begin{cases}x_k,&k\geqslant1,\\0,&k\leqslant0.\end{cases}
\]
Его образ обозначим
\[
  X_+=\operatorname{Im}P_+
  =\{x\in\ell^1(\ZZ):x_k=0\text{ при }k\leqslant0\}.
\]
Это замкнутое подпространство, естественно отождествляемое с
$\ell^1(\mathbb N)$, где $\mathbb N=\{1,2,\ldots\}$.
Положим также $P_-=I-P_+$ и $X_-=\operatorname{Im}P_-$.

Пусть $A$ --- двусторонний оператор свёртки с ядром
$(a_k)\in\ell^1(\ZZ)$ и символом $a(t)=\sum_{k\in\ZZ}a_kt^k\in W$:
\[
  (Ax)_k=\sum_{j\in\ZZ}a_{k-j}x_j.
\]

\begin{definition}
\emph{Оператором Винера--Хопфа} с символом $a$ называется сжатие
\[
  \boxed{T_a=(P_+AP_+)\big|_{X_+}\colon X_+\longrightarrow X_+}.
\]
\end{definition}

Для $x\in X_+$ имеем $P_+x=x$, поэтому
\[
  (T_ax)_k=\sum_{j=1}^{\infty}a_{k-j}x_j,\qquad k\geqslant1.
\]
Оператор ограничен, причём $\norm{T_a}\leqslant\norm{A}
=\sum_k|a_k|$. В стандартном базисе $X_+$ его матрица равна
\[
  T_a=\begin{pmatrix}
    a_0&a_{-1}&a_{-2}&a_{-3}&\cdots\\
    a_1&a_0&a_{-1}&a_{-2}&\cdots\\
    a_2&a_1&a_0&a_{-1}&\cdots\\
    a_3&a_2&a_1&a_0&\cdots\\
    \vdots&\vdots&\vdots&\vdots&\ddots
  \end{pmatrix}.
\]
Это односторонняя бесконечная матрица с постоянными диагоналями.

\subsection{Как действуют проекторы на матрицу}

Столбец матрицы оператора с номером $j$ состоит из координат $Ae^{(j)}$.
Умножение справа на $P_+$ обнуляет столбцы с $j\leqslant0$,
а умножение слева --- строки с $k\leqslant0$.
Относительно разложения $X=X_-\oplus X_+$ это выглядит так:
\[
  A=\begin{pmatrix}A_{--}&A_{-+}\\A_{+-}&A_{++}\end{pmatrix},
  \quad
  P_+=\begin{pmatrix}0&0\\0&I_{X_+}\end{pmatrix},
\]
\[
  AP_+=\begin{pmatrix}0&A_{-+}\\0&A_{++}\end{pmatrix},\qquad
  P_+AP_+=\begin{pmatrix}0&0\\0&A_{++}\end{pmatrix}.
\]
На $X_+$ оператор $T_a$ представлен блоком $A_{++}$.

\begin{remark}
В рукописи обозначение $P_+AP_+$ используется и для оператора на $X$,
и для его ограничения на $X_+$. Здесь эти роли различены.
На всём $X$ такой оператор обнуляет $X_-$ и потому не может быть
обратимым. Все утверждения об обратимости сжатий ниже относятся
к $X_+$; на нём $P_+$ играет роль единичного оператора.
\end{remark}

\section{Треугольные свёртки и проектор}

Обозначим через $\Conv$ алгебру двусторонних операторов свёртки,
а через $\Conv_+$ и $\Conv_-$ её подалгебры с символами из $W_+$
и $W_-$. В рукописи эти алгебры обозначены $\mathcal V,\mathcal V_\pm$;
обозначение $\Conv$ позволяет отличать их от сдвига $V$.

Так как матричный элемент $A_{kj}=a_{k-j}$, получаем
\begin{align*}
  A\in\Conv_+&\ \Longleftrightarrow\ a_n=0\ (n<0)
     \ \Longleftrightarrow\ A\text{ нижнетреуголен},\\
  A\in\Conv_-&\ \Longleftrightarrow\ a_n=0\ (n>0)
     \ \Longleftrightarrow\ A\text{ верхнетреуголен}.
\end{align*}

\begin{lemma}
Для $A\in\Conv$ справедливо
\[
  \boxed{A\in\Conv_+\quad\Longleftrightarrow\quad P_+AP_+=AP_+}.
\]
Это равенство операторов на всём $X$.
\end{lemma}

\begin{proof}
Если $a_n=0$ при $n<0$, то для $k\leqslant0$ и $j\geqslant1$
имеем $k-j<0$ и $a_{k-j}=0$. Поэтому $A$ переводит $X_+$ в себя,
и дополнительное применение $P_+$ слева не меняет $AP_+$.

Обратно, применим равенство к $e^{(1)}$. Так как
$Ae^{(1)}\in X_+$, для всех $k\leqslant0$ выполнено
\[
  0=(Ae^{(1)})_k=a_{k-1}.
\]
Индексы $k-1$ пробегают все отрицательные целые, следовательно,
$A\in\Conv_+$.
\end{proof}

\begin{lemma}
Для $A\in\Conv$ справедливо
\[
  \boxed{A\in\Conv_-\quad\Longleftrightarrow\quad P_+AP_+=P_+A}.
\]
Это также равенство операторов на всём $X$.
\end{lemma}

\begin{proof}
Если $a_n=0$ при $n>0$, то для $k\geqslant1$ и $j\leqslant0$
имеем $a_{k-j}=0$. Значит, $P_+AP_-=0$ и
$P_+A=P_+A(P_++P_-)=P_+AP_+$.

Обратно, применим равенство к $e^{(0)}$. Левая часть равна нулю,
поэтому для всех $k\geqslant1$
\[
  0=(P_+Ae^{(0)})_k=a_k.
\]
Следовательно, $A\in\Conv_-$.
\end{proof}

В рукописи первая лемма доказана матричными схемами, а вторая
предложена для самостоятельной проверки. Координатные доказательства
выше разворачивают обе схемы при принятой границе $k=1$.

\section{Факторизация оператора Винера--Хопфа}

Пусть $a(t)\ne0$ на $\Gamma$. По общей теореме предыдущей лекции
\[
  a(t)=a_-(t)t^\kappa a_+(t),\qquad
  a_\pm,a_\pm^{-1}\in W_\pm,\qquad \kappa=\ind a.
\]
Соответствующий двусторонний оператор факторизуется как
\[
  A=A_-V^\kappa A_+,
\]
где $(Vx)_k=x_{k-1}$ и символ $V$ равен $t$.
Обозначим на $X_+$
\[
  T_\pm=(P_+A_\pm P_+)\big|_{X_+},\qquad
  S_\kappa=(P_+V^\kappa P_+)\big|_{X_+}.
\]

\begin{theorem}
При указанных условиях
\[
  \boxed{T_a=T_-S_\kappa T_+}.
\]
\end{theorem}

\begin{proof}
По двум леммам $P_+A_-=P_+A_-P_+$ и $A_+P_+=P_+A_+P_+$.
Поэтому на всём $X$ имеем
\begin{align*}
  P_+AP_+
  &=P_+A_-V^\kappa A_+P_+\\
  &=P_+A_-P_+V^\kappa P_+A_+P_+\\
  &=(P_+A_-P_+)(P_+V^\kappa P_+)(P_+A_+P_+),
\end{align*}
где использовано $P_+^2=P_+$. Ограничение на $X_+$ даёт формулу теоремы.
\end{proof}

\begin{remark}
Вставка проекторов обоснована треугольностью крайних множителей.
Для произвольных операторов сжатие произведения не обязано равняться
произведению сжатий: промежуточное проектирование может удалить
координаты, которые следующий оператор вернул бы в $X_+$.
\end{remark}

\section{Обратимость крайних множителей}

\begin{theorem}
Операторы $T_-$ и $T_+$ обратимы на $X_+$, причём
\[
  \boxed{T_\pm^{-1}=(P_+A_\pm^{-1}P_+)\big|_{X_+}}.
\]
\end{theorem}

\begin{proof}
Сначала рассмотрим знак минус. Операторы $A_-$ и $A_-^{-1}$
принадлежат $\Conv_-$. Введём $U_-=(P_+A_-^{-1}P_+)|_{X_+}$.
Используя вторую лемму, получаем
\begin{align*}
  (P_+A_-P_+)(P_+A_-^{-1}P_+)
  &=P_+A_-P_+A_-^{-1}P_+\\
  &=P_+A_-A_-^{-1}P_+=P_+.
\end{align*}
При перестановке $A_-$ и $A_-^{-1}$ та же лемма даёт произведение
в обратном порядке, также равное $P_+$. На $X_+$ это означает
$T_-U_-=U_-T_-=I_{X_+}$.

Для знака плюс положим $U_+=(P_+A_+^{-1}P_+)|_{X_+}$.
Теперь оба оператора принадлежат $\Conv_+$. По первой лемме
$P_+A_+^{-1}P_+=A_+^{-1}P_+$, поэтому
\[
  (P_+A_+P_+)(P_+A_+^{-1}P_+)
  =P_+A_+A_+^{-1}P_+=P_+.
\]
Аналогично произведение в обратном порядке равно $P_+$.
Таким образом, $T_+U_+=U_+T_+=I_{X_+}$.
\end{proof}

В рукописи подробно проверен случай минуса, а случай плюса оставлен
для самостоятельного доказательства. Здесь записаны оба.

Факторизация сводит изучение $T_a$ к центральному множителю
$S_\kappa$, поскольку крайние множители обратимы. Это мотивирует
следующий раздел об односторонних обратных операторах.

\section{Левый и правый обратные операторы}

Теперь $X$ --- произвольное банахово пространство,
$A\in\mathcal B(X)$ --- ограниченный линейный оператор.

\begin{definition}
Оператор $A$ называется \emph{обратимым слева}, если существует
ограниченный оператор $C$ такой, что $CA=I$.
Оператор $C$ называют левым обратным к $A$.

Оператор $A$ называется \emph{обратимым справа}, если существует
ограниченный оператор $D$ такой, что $AD=I$.
Оператор $D$ называют правым обратным к $A$.
\end{definition}

Односторонний обратный не обязан быть единственным. Слова
\emph{только слева} означают наличие левого обратного и отсутствие
правого; \emph{только справа} имеют симметричный смысл.

Если существуют оба обратных, они совпадают:
\[
  C=C(AD)=(CA)D=D.
\]
В этом случае $A$ обратим в обычном двустороннем смысле.
Кроме того, наличие левого обратного влечёт инъективность:
$Ax=0\Rightarrow x=CAx=0$. Наличие правого обратного влечёт
сюръективность: $y=A(Dy)$ для каждого $y\in X$.
Эти следствия добавлены для пояснения определений.

\subsection{Пример: два односторонних сдвига}

\emph{Дополнительный пример.} На $\ell^1(\mathbb N)$ положим
\[
  R(x_1,x_2,\ldots)=(0,x_1,x_2,\ldots),\qquad
  L(x_1,x_2,\ldots)=(x_2,x_3,\ldots).
\]
Оба оператора ограничены и имеют норму $1$. Прямое вычисление даёт
\[
  LR=I,\qquad RL=I-E_1,
\]
где $E_1x=(x_1,0,0,\ldots)$.
Следовательно, $R$ имеет левый обратный $L$, но не имеет правого:
его образ состоит из последовательностей с нулевой первой координатой.
Оператор $L$ имеет правый обратный $R$, но не имеет левого:
$Le^{(1)}=0$.

Этот пример связан с центральным множителем факторизации:
\[
  S_1=R,\qquad S_{-1}=L.
\]
В частности, $S_1S_{-1}=I-E_1$, хотя двусторонние сдвиги удовлетворяют
$VV^{-1}=I$. Это конкретный пример того, что сжатие не сохраняет
произведение произвольных операторов.

\section{Степени и произведения односторонне обратимых операторов}

\subsection{Возведение в натуральную степень}

\begin{proposition}
Если $A$ обратим только слева, то $A^n$ обратим только слева
для каждого $n\geqslant1$. Аналогично, если $A$ обратим только
справа, то $A^n$ обратим только справа.
\end{proposition}

\begin{proof}
Пусть $CA=I$. Последовательное сокращение соседних множителей даёт
\[
  C^nA^n=I.
\]
Коммутативность не требуется: например,
$C^2A^2=C(CA)A=CA=I$.
Предположим, что у $A^n$ существует правый обратный $D$.
Тогда
\[
  A^nD=I\quad\Longrightarrow\quad A(A^{n-1}D)=I.
\]
Это даёт правый обратный к $A$, вопреки условию. Значит,
$A^n$ обратим только слева.

Если $AD=I$, то аналогично $A^nD^n=I$. Если бы у $A^n$ был
левый обратный $C$, то
\[
  CA^n=I\quad\Longrightarrow\quad(CA^{n-1})A=I,
\]
и у $A$ существовал бы левый обратный. Это доказывает второй случай.
\end{proof}

\subsection{Умножение на обратимый оператор}

\begin{proposition}
Пусть $B$ двусторонне обратим. Если $A$ обратим только слева,
то $AB$ и $BA$ также обратимы только слева.
Если $A$ обратим только справа, то $AB$ и $BA$ обратимы только справа.
\end{proposition}

\begin{proof}
Сначала пусть $CA=I$. Тогда левые обратные строятся явно:
\[
  (B^{-1}C)(AB)=I,\qquad (CB^{-1})(BA)=I.
\]
Если бы $(AB)D=I$, то $A(BD)=I$; если бы $(BA)D=I$,
то после умножения слева на $B^{-1}$ и справа на $B$ получили бы
$A(DB)=I$. В обоих случаях возник бы правый обратный к $A$.
Следовательно, оба произведения обратимы только слева.

Теперь пусть $AD=I$. Правые обратные задаются формулами
\[
  (AB)(B^{-1}D)=I,\qquad (BA)(DB^{-1})=I.
\]
Если $C(AB)=I$, то $(BC)A=I$; если $C(BA)=I$, то $(CB)A=I$.
Обе возможности противоречат отсутствию левого обратного к $A$.
Следовательно, оба произведения обратимы только справа.
\end{proof}

В рукописи часть симметричных случаев предложена для самостоятельной
проверки; здесь они включены в доказательство. Все формулы сохраняют
порядок множителей.

\section{Устойчивость односторонней обратимости}

Как и обычная обратимость, обратимость только с одной стороны
сохраняется при достаточно малом возмущении. Используем ряд Неймана:
если $\norm{H}<1$, то $I-H$ обратим и
\[
  \norm{(I-H)^{-1}}\leqslant\frac1{1-\norm{H}}.
\]

\subsection{Возмущение оператора, обратимого только слева}

\begin{theorem}
Пусть $A$ обратим только слева и $CA=I$. Если
\[
  \norm{B-A}<\frac1{\norm{C}},
\]
то $B$ также обратим только слева.
\end{theorem}

\begin{proof}
Положим $H=(A-B)C$. Тогда $\norm{H}\leqslant\norm{A-B}\norm{C}<1$,
а значит, $F=I-H$ обратим. При этом
\[
  FA=(I-(A-B)C)A=A-(A-B)CA=B.
\]
Таким образом, $B=FA$ есть произведение обратимого оператора $F$
и оператора $A$, обратимого только слева. По предыдущему утверждению
$B$ обратим только слева.

Можно также явно выписать левый обратный:
\[
  C_B=CF^{-1},\qquad C_BB=CF^{-1}FA=CA=I.
\]
Если бы $BD=I$, то из $FAD=I$ следовало бы $A(DF)=I$,
что противоречит отсутствию правого обратного к $A$.
\end{proof}

В частности, можно выбрать любое $\delta$ с
$0<\delta<1/\norm{C}$ и рассматривать все $B$ с
$\norm{B-A}<\delta$, как это сделано в рукописи.
Дополнительная оценка построенного обратного:
\[
  \norm{C_B}\leqslant
  \frac{\norm{C}}{1-\norm{C}\norm{B-A}}.
\]

\subsection{Возмущение оператора, обратимого только справа}

\begin{theorem}
Пусть $A$ обратим только справа и $AD=I$. Если
\[
  \norm{B-A}<\frac1{\norm{D}},
\]
то $B$ также обратим только справа.
\end{theorem}

\begin{proof}
Этот случай в рукописи оставлен для самостоятельного доказательства.
Положим $H=D(A-B)$ и $F=I-H$. Тогда $\norm{H}<1$, так что $F$ обратим.
Используя $AD=I$, получаем
\[
  AF=A(I-D(A-B))=A-AD(A-B)=B.
\]
Следовательно, $B=AF$ обратим только справа по утверждению о
произведениях. Правый обратный имеет вид
\[
  D_B=F^{-1}D,\qquad BD_B=AFF^{-1}D=AD=I.
\]
Если бы $CB=I$, то $(FC)A=I$, и у $A$ был бы левый обратный,
что невозможно.
\end{proof}

Аналогично справедлива дополнительная оценка
\[
  \norm{D_B}\leqslant
  \frac{\norm{D}}{1-\norm{D}\norm{B-A}}.
\]

\begin{remark}
Здесь утверждается сохранение именно типа <<только слева>> или
<<только справа>>, а не одного равенства с обратным оператором.
Это обеспечивается представлением возмущённого оператора как
произведения исходного с двусторонне обратимым множителем.
\end{remark}

Последняя строка рукописи возвращается к оператору $P_+V^\kappa P_+$.
Его дальнейший разбор на предоставленном скане отсутствует.
Примеры $\kappa=1$ и $\kappa=-1$ выше добавлены для пояснения
односторонней обратимости; общая классификация центрального
множителя здесь не приписывается рукописи.

\end{document}
